problem:  
Let \((M,\mathcal{D},g)\) be a **compact, equiregular, step‑2 sub‑Riemannian manifold** with canonical Popp measure \(\mu\).  
For each \(x\in M\), let \(\exp_x:T_x^*M\to M\) denote the sub‑Riemannian exponential map generated by the Hamiltonian \(H=\tfrac12 g^{-1}\).  For a unit covector \(\lambda\in T_x^*M\), define the **horizontal Jacobian density**
\[
J(x,\lambda,t)\;=\;\det_{\mathcal D}\!\big(D_\lambda\exp_x(t\lambda)\big)
\]
computed using any smooth adapted horizontal frame whose volume form agrees with \(\mu\).  

In contact manifolds, the Agrachev–Barilari–Rizzi (ABR) framework shows that as \(t\to0\),
\[
-\log\frac{J(x,\lambda,t)}{t^{Q-1}} \;=\; \tfrac16\,\mathrm{Ric}_{\mathrm{ABR}}(\lambda)\,t^2+O(t^3),
\qquad Q=\dim_H(M),
\]
so that the quadratic coefficient recovers the canonical Ricci curvature.  

**Problem:**  
Assume now that \((M,\mathcal{D},g)\) is *non‑contact* but still step‑2 and equiregular, possibly admitting abnormal minimizers.  
Investigate whether the second‑order coefficient in the small‑time asymptotic expansion  
\[
-\log\frac{J(x,\lambda,t)}{t^{Q-1}}
     = A(x,\lambda)\,t^2 + o(t^2), \quad t\to0,
\]
exists for all normal covectors \(\lambda\), and if so:

1. Does \(A(x,\lambda)\) define a smooth quadratic form in \(\lambda\) on the horizontal cotangent space, independent of the chosen adapted frame?  
2. If it exists and is intrinsic, does \(A(x,\lambda)\) coincide with \(\mathrm{Ric}_{\mathrm{ABR}}(\lambda)\) in the contact case and extend it naturally in the general step‑2 case?  
3. Can the sign of \(A(x,\lambda)\) be interpreted as the infinitesimal condition guaranteeing **entropy convexity** of measures transported along the normal displacement interpolations generated by the Hamiltonian flow?

In other words, does the \(t^2\)-coefficient of the horizontal Jacobian distortion always define an **intrinsic effective Ricci curvature tensor** governing the analytic behavior of optimal transport near normal geodesics in any step‑2 sub‑Riemannian manifold?

Why is it a "good" problem:  
- **Precision and focus:** This formulation isolates a single invariance and existence question—whether a well‑defined curvature coefficient \(A(x,\lambda)\) exists intrinsically from the small‑time expansion of Jacobian distortion along Hamiltonian geodesics. It avoids broad programmatic claims and centers on a concrete, verifiable geometric object.  
- **Unknown yet accessible:** Even for the simplest non‑contact step‑2 examples (such as the Martinet distribution), this coefficient is not rigorously understood. Determining its existence or invariance would constitute a genuine advance.  
- **Conceptual depth:** The question probes whether the infinitesimal structure underlying synthetic curvature and entropy convexity (known in the Riemannian/contact cases) survives in the presence of abnormal geodesics and bracket degeneracy.  
- **Bridge between fields:** It unites geometric control theory (Jacobian of the exponential map), analytic curvature notions (ABR and entropy convexity), and optimal transport theory—potentially revealing how curvature arises from measure distortion in nonholonomic geometry.  
- **Elegant simplicity:** Expressed entirely through the small‑time behavior of a single explicitly definable quantity, the problem is conceptually transparent but technically profound—precisely the “narrow entrance, profound depth” hallmark of a good research problem.